% Introduction — repositioned as systematic empirical comparison

Variational quantum circuits (VQCs) have emerged as one of the most promising
paradigms for quantum machine learning on near-term intermediate-scale quantum
(NISQ) devices~\cite{cerezo2021variational,benedetti2019parameterized}.
By parameterizing quantum gates and optimizing them via classical gradient
methods, VQCs combine the expressiveness of quantum Hilbert space with the
practicality of hybrid quantum--classical training loops.

Despite growing interest, the design space of VQC architectures remains
poorly understood.  Classical deep learning has benefited enormously from
systematic architectural comparisons---fully-connected networks versus
convolutional networks versus transformers---which yielded clear guidelines on
when each topology excels.  No analogous study exists for VQCs.  Most prior
work proposes a single new circuit ansatz and evaluates it on one or two
benchmarks, making it difficult to draw general conclusions.

In this work, we address this gap with a \emph{systematic empirical
comparison} of four VQC architecture families on tabular regression and
classification tasks:

\begin{enumerate}
  \item \textbf{Multi-layer fully-connected VQC (FC-VQC)} ---
    cascaded VQC blocks with all-to-all inter-block connectivity
    and configurable qubit layouts (e.g.\ $16 \to 4 \to 1$).
  \item \textbf{Residual VQC (ResNet-VQC)} ---
    VQC layers augmented with classical residual (skip) connections,
    inspired by He et al.~\cite{he2016deep}.
  \item \textbf{Quantum Transformer VQC (QT)} ---
    a hybrid architecture where classical self-attention operates on
    quantum-encoded features (Route~A).
  \item \textbf{Fully Quantum Transformer VQC (FQT)} ---
    an architecture where both the attention mechanism and the
    feed-forward network are implemented as parameterized quantum
    circuits (Route~B).
\end{enumerate}

All architectures are evaluated on five datasets (three regression, two
classification) against strong classical baselines including XGBoost,
CatBoost, and parameter-matched MLPs.  We further conduct ablation studies
on the transformer components and measure circuit expressibility using the
Sim et al.~\cite{sim2019expressibility} KL-divergence metric.

Our main contributions are:
\begin{itemize}
  \item A unified experimental framework enabling fair, controlled comparison
    across VQC architectures on identical data splits, training protocols,
    and evaluation metrics.
  \item Evidence that \textbf{FC-VQCs are the most parameter-efficient}
    architecture, achieving regression accuracy competitive with
    gradient-boosted trees using ${\sim}$700 parameters versus
    ${\sim}$13{,}000 for equivalently accurate MLPs.
  \item A \textbf{negative result on quantum attention}: self-attention
    does not consistently help on small tabular data, and removing it can
    improve performance---an important finding for practitioners.
  \item Practical guidance: expressibility saturates at depth~${\approx}\,3$,
    LayerNorm benefits fully quantum architectures on classification, and
    residual connections offer a reliable middle ground between FC-VQC and
    transformer complexity.
\end{itemize}
